\begin{frame}{Explotación \hfill{\tiny\color{TextoSuave}\textbf{WT}: enzima original \quad | \quad \textbf{ENG}: enzima modificada por ingeniería}}



{\color{Cian}\bfseries Resumen de la Sección}\par
\vspace{0.07cm}

\begin{columns}[T,onlytextwidth]
\column{0.47\textwidth}
\begin{block}{Qué Hacer}
\small
\begin{enumerate}
\item \textbf{Glucósidos:} 60 $^{\circ}$C, 40 mM $\rightarrow$ producto útil + selectividad.
\item \textbf{Biomasa:} 70 $^{\circ}$C, 80 mM $\rightarrow$ explotar ventana térmica.
\item \textbf{Bagazo:} 60 $^{\circ}$C, 40 mM $\rightarrow$ estabilidad + tolerancia.
\item \textbf{Alimentos:} 60 $^{\circ}$C, 40 mM $\rightarrow$ vida operativa.
\end{enumerate}
\end{block}

\column{0.49\textwidth}
\begin{alertblock}{Prioridad Productiva}
\centering\small\bfseries
Glucósidos 5.95$\times$\\
$>$ Biomasa 4.88$\times$\\
$>$ Bagazo 4.33$\times$\\
$>$ Alimentos 2.78$\times$
\end{alertblock}

\vspace{0.08cm}
\begin{beamercolorbox}[wd=\linewidth,sep=0.10cm,rounded=true]{block body}
\scriptsize
\textbf{Por qué:} glucósidos combina IC = 2.56 con selectividad/tolerancia; FP = 1.87 amplifica la mejora molecular.
\end{beamercolorbox}
\end{columns}

\vspace{0.08cm}
\begin{block}{Lectura Económica}
\centering\small
Con los datos actuales se identifica \textbf{potencial productivo}. La rentabilidad neta se confirma con precio, costo ENG y energía.
\end{block}

\vfill
\begin{columns}[b,onlytextwidth]
\column{0.23\textwidth}
\boton{resultados}{Anterior: Resultados}
\column{0.52\textwidth}
\column{0.23\textwidth}
\raggedleft\boton{conclusion}{Siguiente: Conclusión}
\end{columns}
\end{frame}
